% Modelo de LaTeX para as listas de MAC0320/5771.
%
% Podem utiliza-lo como base para a solução de vocês, fazendo as alterações
% necessárias.

\documentclass[11pt]{amsart}

% Muitos pacotes úteis (nem todos são utilizados aqui)
\usepackage[utf8]{inputenc}
\usepackage[T1]{fontenc}
\usepackage[brazilian]{babel}
\usepackage{calc}
\usepackage{xargs}
\usepackage{ifthen}
\usepackage{xspace}
\usepackage{xkeyval}
\usepackage{blindtext}
\usepackage{amsfonts}
\usepackage{amsmath}
\usepackage{amssymb}
\usepackage{amsthm}
\usepackage{geometry}
\usepackage{booktabs}
\usepackage{hyperref}
\usepackage{cleveref}
\usepackage{setspace}
\usepackage{csquotes}
\usepackage{enumitem}
\usepackage{multicol}
\usepackage{caption}
\usepackage{subcaption}
\usepackage{tcolorbox}
\usepackage{url}
\usepackage{tikz}

% Definindo algumas coisas importantes
\theoremstyle{plain}
\newtheorem{theorem}{Teorema}
\theoremstyle{definition}
\newtheorem{exercise}{Exercício}

% Informações básicas
\title{Lista 2 - MAC0320/MAC5770 - Árvores e Grafos Eulerianos}
\author{Fábio Botler \and Bruno Akamine}

\begin{document}

\maketitle

\begin{center}
  \begin{tcolorbox}[width=0.65\textwidth,colframe=red,colback=red!20]
    \textsc{Entrega até o dia 16/4 às 23:59 os seguintes exercícios: 2, 3, 4, 7, 9}
  \end{tcolorbox}
\end{center}

\vspace{1cm}

\textbf{Orientações (caso não sejam cumpridas, poderá haver desconto de nota):}
\begin{enumerate}
  \item Resolva um exercício por página.
  \item Entregue apenas os exercícios solicitados.
  \item A lista deve estar legível.
  \item Envie um único arquivo PDF. O nome do arquivo deve seguir o formato NUSP.pdf.
  \item Cada aluno tem um crédito de atraso de 7 dias. Isso é, a soma total dos
    atrasos de todas as listas não deve ultrapassar 7 dias. Listas entregues
    após o prazo não serão aceitas caso o aluno já tenha extrapolado os créditos
\end{enumerate}
\textbf{Sugestão:} use o chatGPT para passar a solução de cada exercício para \LaTeX, mas não esqueça de verificar o trabalho feito por ele.

\vspace{5mm}

\begin{exercise}
  Prove que se $G$ é uma árvore tal que $\Delta(G) \ge k$, então $G$ tem pelo
  menos $k$ folhas.
\end{exercise}

\begin{exercise}
  Prove que todo grafo conexo $G$, simples e não-trivial, tem uma árvore
  geradora $T$, tal que $G-E(T)$ é desconexo.
\end{exercise}

\begin{exercise}
  Seja $G$ um grafo conexo, $T_{1}, T_{2}$ árvores geradoras distintas de $G$, e
  seja, $e_{1}$ uma aresta de $T_{1}$.
  Prove que existe uma aresta $e_{2}$ em $T_{2}$ tal que $T_{1} - e_{1} +
  e_{2}$ é uma árvore geradora de $G$.
\end{exercise}

\begin{exercise}
  Uma \textit{$k$-coloração} dos vértices de um grafo é uma atribuição de no máximo $k$
  cores distintas aos vértices desse grafo, tal que vértices adjacentes recebem
  cores distintas. Dizemos que um grafo é \textit{equi-bicolorido} se tem uma
  $2$-coloração com igual número de vértices de cada cor.
  Prove que toda árvore equi-bicolorida tem pelo menos uma folha de cada cor.
\end{exercise}

\begin{exercise}
  De uma segunda prova para o exercício anterior.
\end{exercise}

\begin{exercise}
  Seja $\mathcal{F}(T)$ o número de folhas de uma árvore $T$ e seja $S \subset
  V(T)$ o conjunto dos vértices com grau no mínimo 3. Mostre que se $T$
  é uma árvore não-trivial, então \[
    \mathcal{F}(T) = 2 + \sum_{v \in S} \left( d(v) - 2 \right)
  .\]
\end{exercise}

\begin{exercise}
  Prove que um grafo conexo $G$ é euleriano se, e só se, $G$ contém ciclos
  $C_{1},C_{2},\ldots,C_{k}$ dois a dois disjuntos nas arestas tais que $E(G) =
  E(C_{1}) \cup E(C_{2}) \cup \cdots \cup E(C_{k})$.
\end{exercise}

\begin{exercise}
  Prove que se $G$ é um grafo conexo tal que cada uma das suas arestas pertence
  a um número ímpar de ciclos, então $G$ é euleriano.
\end{exercise}

\begin{exercise}
    Seja $G$ um grafo conexo e seja $S\subset V(G)$ um conjunto e vértices tal que $|S|\geq 2$. Prove que se $N_G(v)\setminus S\neq \emptyset$, para cada $v\in S$ e $G[V(G)\setminus S]$ é conexo, então $G$ tem uma árvore $T$ tal que o conjunto de folhas de $T$ é exatamente $S$.
\end{exercise}
\end{document}
